% -----------------------------------------------------------
%  Anexo — Definiciones y Abreviaciones del Negocio
% -----------------------------------------------------------

Este anexo recopila los términos y conceptos del dominio de administración de comunidades residenciales utilizados a lo largo del informe.

\begin{description}
    \item[Administrador de comunidad:] Persona natural o jurídica designada por la asamblea de copropietarios para gestionar la operación, finanzas y mantención de un edificio o condominio, conforme a la Ley N\textsuperscript{o}~21.442.

    \item[Alícuota:] Porcentaje de participación de cada unidad habitacional en los gastos comunes del edificio, calculado proporcionalmente según la superficie o criterios definidos en el reglamento de copropiedad.

    \item[Área común:] Espacios compartidos por los copropietarios, tales como quinchos, piscinas, salones multiuso, estacionamientos de visitas y áreas verdes.

    \item[Comité de administración:] Órgano elegido por la asamblea de copropietarios que fiscaliza la gestión del administrador y toma decisiones en representación de la comunidad.

    \item[Copropiedad inmobiliaria:] Régimen jurídico en el que un inmueble es compartido por múltiples propietarios, cada uno titular de una unidad privada y copropietario de los bienes comunes.

    \item[Encomienda:] Paquete o envío recibido en la conserjería del edificio a nombre de un residente, sujeto a registro, custodia y entrega controlada.

    \item[Gastos comunes:] Cobro mensual que cada unidad habitacional debe pagar para cubrir los costos de administración, mantención, servicios básicos y reparaciones de los bienes comunes del edificio.

    \item[Marcha blanca:] Período de transición durante el cual un nuevo sistema opera en paralelo con los procesos manuales existentes, permitiendo validar su funcionamiento sin interrumpir la operación.

    \item[Morosidad:] Estado de deuda de un copropietario que no ha pagado los gastos comunes dentro del plazo establecido. La Ley 21.442 permite restricciones al moroso (bloqueo de reservas, entre otras).

    \item[Pre-autorización:] Permiso otorgado por un residente a un visitante específico para ingresar al edificio en una fecha y horario determinados, verificable mediante código QR temporal.

    \item[Proptech:] Sector de la industria tecnológica que desarrolla soluciones digitales para el mercado inmobiliario y la gestión de propiedades.

    \item[SaaS:] \textit{Software as a Service}. Modelo de distribución de software en el que la aplicación se ofrece como servicio en la nube mediante suscripción mensual.

    \item[Unidad habitacional:] Departamento, casa o local individual dentro de una copropiedad inmobiliaria, identificado por un número o código único.
\end{description}
